\subsection{数据范围与样本选择}\label{sec:data}
\subsubsection{完整股票范围与筛选规则}
初始分钟资料覆盖2017年1月3日至2019年12月31日的731个沪市交易日，共1573只股票、
996\,288个有记录的“股票—交易日”。2017、2018和2019年分别有1387、1447和1568只股票。
各股票的可用日期不同，1573只股票并非每天都有记录；逐股覆盖与筛选结果见补充材料。

筛选只使用前584个名义训练日，即2017年1月3日至2019年5月29日。
有记录日比例不低于80\%的股票有1225只，将门槛提高到95\%后剩952只；
再要求完整日比例不低于90\%、日成交额中位数为正，得到158只候选。
按训练期日成交额中位数降序取前8只，中国平安排名第一，作为主例。
缺记录日的成交额按零计入584日的中位数。80\%初筛不会排除满足最终95\%门槛的股票。

图\ref{fig:universe}(a)每个点代表一只股票。横轴为“有记录日数$/584\times100\%$”，
衡量缺了多少交易日；纵轴为“完整日数$/$有记录日数$\times100\%$”，衡量日内分钟记录是否齐全。
没有训练期记录的股票置于原点，仅用来标记缺失。完整日要求09:30--11:29、
13:00--14:56共237个核心分钟标签齐备，且15:00收盘记录存在。
这一口径考虑了收盘前部分分钟标签的系统性缺失，不能简单以每天的总行数代替。
上述统计共覆盖774\,275个训练期股票日、175\,980\,062条分钟记录。

\begin{figure}[!htbp]\centering\begin{tikzpicture}
\begin{groupplot}[group style={group size=2 by 1,horizontal sep=1.5cm},width=.455\linewidth,height=6.2cm,
  grid=major,grid style={gray!15},tick label style={font=\small},label style={font=\small},title style={font=\small}]
\nextgroupplot[xmin=-2,xmax=102,ymin=-2,ymax=102,xlabel={训练期有记录日比例（\%）},ylabel={有记录日中的完整日比例（\%）},title={(a) 全部 1573 只股票}]
\addplot[only marks,mark=*,mark size=.65pt,gray!55] table[x=presence_pct,y=complete_pct,col sep=comma]{universe_all.csv};
\addplot[only marks,mark=*,mark size=1pt,blue!65] table[x=presence_pct,y=complete_pct,col sep=comma]{universe_eligible.csv};
\addplot[only marks,mark=star,mark size=2pt,orange!90!black] table[x=presence_pct,y=complete_pct,col sep=comma]{universe_selected.csv};
\draw[dashed,black!65] (axis cs:95,0)--(axis cs:95,100);
\draw[dashed,black!65] (axis cs:0,90)--(axis cs:100,90);
\nextgroupplot[xmin=1,xmax=158,ymode=log,xlabel={合格候选的成交额排名},ylabel={日成交额中位数（亿元）},title={(b) 158 只合格候选的排序}]
\addplot[thick,blue!65,no marks] table[x=eligible_rank,y=amount_yi,col sep=comma]{universe_eligible.csv};
\addplot[only marks,mark=*,mark size=1.7pt,orange!90!black] table[x=eligible_rank,y=amount_yi,col sep=comma]{universe_selected.csv};
\draw[dashed,black!65] (axis cs:8,.01)--(axis cs:8,100);
\end{groupplot}
\end{tikzpicture}

\caption{完整股票范围与样本筛选。(a) 灰点为全部1573只股票，蓝点为158只合格候选，橙色为最终8只；重合点可能遮挡。
虚线为95\%覆盖率和90\%完整日门槛。(b) 合格候选按训练期成交额中位数排序，橙色为前8名，纵轴取对数。}
\label{fig:universe}\end{figure}
\begin{table}[!htbp]\centering\small
\caption{入选的八只股票。排序只依据训练期日成交额中位数；完整日比例的分母为有记录日数，覆盖率的分母为584个名义训练日。}
\label{tab:selected_stocks}
\begin{tabular}{rlcrrrr}\toprule
排名 & 股票 & 证券代码 & 有记录日 & 覆盖率(\%) & 完整日(\%) & 成交额(亿元)\\\midrule
1 & 中国平安 & 601318.SH & 584 & 100.00 & 100.00 & 39.17\\
2 & 贵州茅台 & 600519.SH & 584 & 100.00 & 100.00 & 22.55\\
3 & 中信证券 & 600030.SH & 584 & 100.00 & 98.29 & 15.62\\
4 & 伊利股份 & 600887.SH & 582 & 99.66 & 99.66 & 13.46\\
5 & 招商银行 & 600036.SH & 584 & 100.00 & 100.00 & 12.44\\
6 & 华友钴业 & 603799.SH & 583 & 99.83 & 98.11 & 12.35\\
7 & 方大炭素 & 600516.SH & 583 & 99.83 & 99.49 & 11.01\\
8 & 中国建筑 & 601668.SH & 584 & 100.00 & 100.00 & 10.75\\
\bottomrule\end{tabular}\end{table}

\FloatBarrier

\subsubsection{样本划分与观测特征}
731个交易日按时间顺序以80\%、10\%、10\%划分训练、验证与测试段，
按累计比例向下取整得到584、73、74日。沿用每段末尾剔除5个交易日的规则，
实际保留579、68、69日，各股票本身缺失的日期不补值。保留训练段为2017年1月3日至2019年5月22日，
验证段为2019年5月30日至9月3日，测试段为2019年9月11日至12月24日。
选股使用584个名义训练日，归一化尺度、观测窗口及模型参数只用保留的训练收益估计。

价格按上午、下午分别聚合为5分钟区间，每个区间要求至少80\%的必需分钟标签，
15:00记录并入末个区间。只计算同一时段内相邻区间的收盘收益，完整日得到46个收益。
样本量门槛中的完整日按46条有效5分钟收益计数，与选股阶段的分钟标签完整性分别核对。
新验证、测试段通常只有3128和3174条收益，均不足既定的“60个完整日且5000条收益”门槛，
以下将其作为短留出段检查。

2026年1月5日至8月31日为后续年份检验。7股各有160日、7360条收益，满足上述门槛；
华友钴业仅有6月30日至8月31日的45日、2070条收益，单列为短样本。
2020年下半年的分钟资料仅覆盖所选八股中的两只，2021--2025年无可用分钟记录，
不将这些年份补齐或拼接成连续样本；日频资料也不替代缺失的分钟数据。
评价采用上述时间划分，检验范围见第\ref{sec:oos_method}节。
后续行情不参与本次选股、归一化常数估计或参数重估。

\begin{table}[!htbp]\centering\small
\caption{八股有效收益样本与训练期特征。各样本列为“收益条数/交易日数”；标准差以原始对数收益的百分数表示，零收益比例未作平滑。}
\label{tab:eight_counts}
\begin{tabular}{lrrrrrr}\toprule
股票 & 训练 & 验证 & 测试 & 2026年 & 标准差(\%) & 零收益(\%)\\\midrule
中国平安 & 26\,634/579 & 3\,128/68 & 3\,174/69 & 7\,360/160 & 0.2211 & 5.32\\
贵州茅台 & 26\,634/579 & 3\,128/68 & 3\,174/69 & 7\,360/160 & 0.2290 & 1.66\\
中信证券 & 26\,174/569 & 3\,128/68 & 3\,128/68 & 7\,360/160 & 0.2498 & 14.95\\
伊利股份 & 26\,450/575 & 3\,128/68 & 3\,174/69 & 7\,360/160 & 0.2695 & 9.21\\
招商银行 & 26\,634/579 & 3\,128/68 & 3\,174/69 & 7\,360/160 & 0.2423 & 9.28\\
华友钴业 & 26\,128/568 & 3\,128/68 & 3\,174/69 & 2\,070/45 & 0.4511 & 3.62\\
方大炭素 & 26\,542/577 & 3\,128/68 & 3\,174/69 & 7\,360/160 & 0.3953 & 10.91\\
中国建筑 & 26\,634/579 & 3\,128/68 & 3\,174/69 & 7\,360/160 & 0.2293 & 29.73\\
\bottomrule\end{tabular}\end{table}

中国平安训练期日成交额中位数最高，且584日均符合完整性条件，适合展示计算过程。
成交活跃度与完整性筛选也使样本偏向交易较活跃的股票；八股比较用于检查个股差异，
不将结果推广为全部A股的共同规律。图\ref{fig:panel_data}展示八股在全部可用日期上的波动变化。
\clearpage
\begin{figure}[p]\centering\begin{tikzpicture}
\begin{groupplot}[group style={group size=2 by 4,horizontal sep=1.7cm,vertical sep=1.25cm},
width=.44\linewidth,height=4.3cm,xmin=0,xmax=929,ymin=0,unbounded coords=jump,
xtick={0,244,487,770},xticklabels={2017,2018,2019,2026},
tick label style={font=\scriptsize},title style={font=\small},label style={font=\small},
grid=major,grid style={gray!15},ylabel={日内标准差（\%）}]
\nextgroupplot[title={中国平安 601318.SH}]
\addplot[blue!75,line width=.45pt,no marks] table[x=x,y=vol_pct,col sep=comma]{panel_vol_601318_old.csv};
\addplot[orange!85!black,line width=.55pt,no marks] table[x=x,y=vol_pct,col sep=comma]{panel_vol_601318_new.csv};
\nextgroupplot[title={贵州茅台 600519.SH}]
\addplot[blue!75,line width=.45pt,no marks] table[x=x,y=vol_pct,col sep=comma]{panel_vol_600519_old.csv};
\addplot[orange!85!black,line width=.55pt,no marks] table[x=x,y=vol_pct,col sep=comma]{panel_vol_600519_new.csv};
\nextgroupplot[title={中信证券 600030.SH}]
\addplot[blue!75,line width=.45pt,no marks] table[x=x,y=vol_pct,col sep=comma]{panel_vol_600030_old.csv};
\addplot[orange!85!black,line width=.55pt,no marks] table[x=x,y=vol_pct,col sep=comma]{panel_vol_600030_new.csv};
\nextgroupplot[title={伊利股份 600887.SH}]
\addplot[blue!75,line width=.45pt,no marks] table[x=x,y=vol_pct,col sep=comma]{panel_vol_600887_old.csv};
\addplot[orange!85!black,line width=.55pt,no marks] table[x=x,y=vol_pct,col sep=comma]{panel_vol_600887_new.csv};
\nextgroupplot[title={招商银行 600036.SH}]
\addplot[blue!75,line width=.45pt,no marks] table[x=x,y=vol_pct,col sep=comma]{panel_vol_600036_old.csv};
\addplot[orange!85!black,line width=.55pt,no marks] table[x=x,y=vol_pct,col sep=comma]{panel_vol_600036_new.csv};
\nextgroupplot[title={华友钴业 603799.SH}]
\addplot[blue!75,line width=.45pt,no marks] table[x=x,y=vol_pct,col sep=comma]{panel_vol_603799_old.csv};
\addplot[orange!85!black,line width=.55pt,no marks] table[x=x,y=vol_pct,col sep=comma]{panel_vol_603799_new.csv};
\nextgroupplot[title={方大炭素 600516.SH}]
\addplot[blue!75,line width=.45pt,no marks] table[x=x,y=vol_pct,col sep=comma]{panel_vol_600516_old.csv};
\addplot[orange!85!black,line width=.55pt,no marks] table[x=x,y=vol_pct,col sep=comma]{panel_vol_600516_new.csv};
\nextgroupplot[title={中国建筑 601668.SH}]
\addplot[blue!75,line width=.45pt,no marks] table[x=x,y=vol_pct,col sep=comma]{panel_vol_601668_old.csv};
\addplot[orange!85!black,line width=.55pt,no marks] table[x=x,y=vol_pct,col sep=comma]{panel_vol_601668_new.csv};
\end{groupplot}\end{tikzpicture}
\caption{八只股票的真实收益波动。每点为当日全部有效5分钟收益的样本标准差。
蓝色为2017--2019年，橙色为2026年；两段之间留空，缺失年份不插值。
横轴按各段交易日顺序排列，压缩两段之间的日历间隔；纵轴独立取值以保留各股变化。
图中包括划分边界处仅用于展示的日期，这些日期不进入训练或留出评分。}
\label{fig:panel_data}\end{figure}\clearpage
